% Part: proof-theory
% Chapter: proof-search
% Section: tableaux

\documentclass[../../../include/open-logic-section]{subfiles}

\begin{document}

\olfileid{pt}{ps}{tab}

\olsection{કોષ્ટક-વૃક્ષો}

કોષ્ટક-વૃક્ષ તંત્રો સિક્વન્ટ-કલન સાથે નજીકનો સંબંધ ધરાવતાં
!!{proof} તંત્રો છે. હાથથી અને સૉફ્ટવેરમાં બંને રીતે
સાબિતી-શોધ માટે તે ખાસ અનુકૂળ છે. કોષ્ટક-વૃક્ષમાં
સિક્વન્ટોનું વૃક્ષ બનાવવાને બદલે !!a{proof} એ
!!{formula}sનું વૃક્ષ છે. જોકે સ્વાભાવિક નિગમનથી
વિપરીત, તેના નિયમો સિક્વન્ટ-કલનના નિયમોને અનુસરે છે.
કોષ્ટક-વૃક્ષની !!{proof} હકીકતમાં નિદર્શની સ્પષ્ટ શોધ
છે. તેથી તેને ક્યારેક \emph{અર્થલક્ષી} કોષ્ટક-વૃક્ષ
અથવા \emph{સત્ય-વૃક્ષ} પણ કહે છે.

\emph{ચિહ્નિત} કોષ્ટક-વૃક્ષ ચિહ્નિત !!{formula}sનું
વૃક્ષ છે; તે $\sFmla{\True}{!A}$ અથવા
$\sFmla{\False}{!A}$ સ્વરૂપનાં હોઈ શકે.
વૃક્ષો નીચેની દિશામાં વધે છે. શરૂઆતમાં જે
!!{formula}s હોય છે તે આપણે શું !!{prove} કરવું છે
તે નક્કી કરે છે. દાખલા તરીકે,
$!A, !B \Sequent !C, !D$ સિક્વન્ટને !!{prove}
કરવા માટે તે
$\sFmla{\True}{!A}, \sFmla{\True}{!B},
\sFmla{\False}{!C}, \sFmla{\False}{!D}$થી
શરૂ થાય.
\sourcecorrection{OLMD-278}{મૂળ ઉદાહરણમાં $!D$ને ભૂલથી સત્ય-ચિહ્ન અપાયું છે; તરત પછીનું ગદ્ય $!C$ અને $!D$ બંને અસત્ય રાખે છે, તેથી અહીં અસત્ય-ચિહ્ન કર્યું છે.}
$!A$ અને $!B$ને સત્ય તથા $!C$ અને $!D$ને
અસત્ય કરતી !!^a{structure} એ
$!A, !B \Sequent !C, !D$ અપ્રામાણ્ય છે તે
દર્શાવતી !!a{structure} છે.

કોષ્ટક-વૃક્ષના પાન, એટલે સૌથી નીચેના શિરોબિંદુ,
શાખા-વિસ્તારના નિયમોથી આગળ વધારી શકાય છે. આ
નિયમો એ જ શાખામાં નવા શિરોબિંદુઓ ઉમેરે છે અથવા
નીચે બે સમકક્ષ શિરોબિંદુઓ ઉમેરી શાખાને બે ભાગે
વહેંચે છે. એક જ શાખામાં પાનની ઉપરના કોઈ પણ
ચિહ્નિત !!{formula} પર નિયમ લાગુ થઈ શકે છે.
શાસ્ત્રીય કોષ્ટક-વૃક્ષ તંત્ર \Log{Tc}ના
શાખા-વિસ્તારના નિયમો \olref{tab:Tc}માં આપ્યા છે.

\subfile{rules-Tc}

જોઈ શકાય છે કે આ નિયમો \Log{G3c}ના નિયમો ઊલટા
મૂકેલા છે. $\True$ અને $\False$ ચિહ્નો સૂચવે છે કે
!!a{formula} સિક્વન્ટની અનુક્રમે ડાબી અને જમણી
બાજુએ મુખ્ય સૂત્ર છે.

કોષ્ટક-વૃક્ષની શાખામાં
$\sFmla{\True}{!A}$ અને $\sFmla{\False}{!A}$
બંને હોય, તો તે \emph{બંધ} છે. દરેક શાખા બંધ
હોય, તો આખું કોષ્ટક-વૃક્ષ બંધ કહેવાય.
દાખલા તરીકે,
$!A \lor (!B \land !C) \Sequent (!A \lor !B) \land
(!A \lor !C)$ માટેનું આ બંધ કોષ્ટક-વૃક્ષ છે;
તે $\sFmla{\True}{!A \lor (!B \land !C)}$
અને $\sFmla{\False}{(!A \lor !B) \land (!A \lor !C)}$થી
શરૂ થાય છે:
\sourcecorrection{OLMD-279}{મૂળમાં clause tableau લખાયું છે, પરંતુ દર્શાવેલા વૃક્ષની બધી શાખાઓ બંધ છે; અહીં બંધ કોષ્ટક-વૃક્ષ કર્યું છે.}
\begin{oltableau}
  [\sFmla{\True}{\formula{A} \lor (\formula{B} \land \formula{C})},
    just=\TAss, checked
    [\sFmla{\False}{(\formula{A} \lor \formula{B}) \land
        (\formula{A} \lor \formula{C})}, just=\TAss, checked
      [\sFmla{\True}{\formula{A}}, just={\TRule{\True}{\lor}[1]}
        [\sFmla{\False}{\formula{A} \lor \formula{B}},
          just={\TRule{\False}{\land}[2]}, checked
          [\sFmla{\False}{\formula{A}},
            just={\TRule{\False}{\lor}[4]}
            [\sFmla{\False}{\formula{B}},
              just={\TRule{\False}{\lor}[4]}, close]
          ]
        ]
        [\sFmla{\False}{\formula{A} \lor \formula{C}},
          just={\TRule{\False}{\land}[2]}, checked
          [\sFmla{\False}{\formula{A}},
            just={\TRule{\False}{\lor}[4]}
            [\sFmla{\False}{\formula{C}},
              just={\TRule{\False}{\lor}[4]}, close]
          ]
        ]
      ]
      [\sFmla{\True}{\formula{B} \land \formula{C}},
        just={\TRule{\True}{\lor}[1]}, checked
        [\sFmla{\False}{\formula{A} \lor \formula{B}},
          just={\TRule{\False}{\land}[2]}, checked
          [\sFmla{\True}{\formula{B}},
            just={\TRule{\True}{\land}[3]}, move by=2
            [\sFmla{\True}{\formula{C}},
              just={\TRule{\True}{\land}[3]}
              [\sFmla{\False}{\formula{A}},
                just={\TRule{\False}{\lor}[4]}
                [\sFmla{\False}{\formula{B}},
                  just={\TRule{\False}{\lor}[4]},close
                ]
              ]
            ]
          ]
        ]
        [\sFmla{\False}{\formula{A} \lor \formula{C}},
          just={\TRule{\False}{\land}[2]}, checked
          [\sFmla{\True}{\formula{B}},
            just={\TRule{\True}{\land}[3]}, move by=2
              [\sFmla{\True}{\formula{C}},
                just={\TRule{\True}{\land}[3]}
                [\sFmla{\False}{\formula{A}},
                  just={\TRule{\False}{\lor}[4]}
                  [\sFmla{\False}{\formula{C}},
                  just={\TRule{\False}{\lor}[4]},close
                ]
              ]
            ]
          ]
        ]
      ]
    ]
  ]
\end{oltableau}

ખરાની નિશાનીઓ સૂચવે છે કે તેમની નીચેની બધી
શાખાઓમાં સંબંધિત !!a{formula} પર નિયમ લાગુ
કરી દેવાયો છે. $\otimes$ ચિહ્ન બંધ શાખા સૂચવે છે.
નિયમોની પાછળની પંક્તિ-સંખ્યાઓ કયા ચિહ્નિત
!!{formula}s પર નિયમ લાગુ કર્યો છે તે બતાવે છે,
અર્થાત્ એ જ શાખામાં દર્શાવેલી પંક્તિ પરના
!!{formula} પર.

કોષ્ટક-વૃક્ષમાં સાબિતી-શોધ સરળ છે. કોષ્ટક-વૃક્ષ
તબક્કાવાર બનાવીએ. તબક્કા~$0$એ શરૂઆતનાં
!!{formula}s સૌથી ઉપર લખો. તબક્કા~$n+1$એ
દરેક પાન માટે આ કરો: જે ચિહ્નિત સૂત્ર પર
હજી ખરાની નિશાની નથી તેમાંથી સૌથી ઉપરનું
પ્રથમ સૂત્ર શોધો. તેના પર શાખા-વિસ્તારનો નિયમ
લાગુ કરો. એ ચિહ્નિત સૂત્ર ધરાવતી દરેક
શાખામાં નિયમ લાગુ થઈ જાય, પછી તેના પર ખરાની
નિશાની કરો. (ઉપરનું ઉદાહરણ આ અલ્ગોરિધમથી
બનાવ્યું છે.)

$\sFmla{\True}{\lforall}$ અને
$\sFmla{\False}{\lexists}$ સ્વરૂપનાં ચિહ્નિત
!!{formula}s પર કદી ખરાની નિશાની થતી નથી,
એટલે તેમની વિશેષ સંભાળ લેવી પડે. તે હાજર હોય,
તો $\TRule{\True}{\forall}$ અને
$\TRule{\False}{\lexists}$ નિયમો દરેક લાગુ પડતા
પદ માટે આખરે લાગુ થાય તેની ખાતરી કરવી પડે.
દાખલા તરીકે, દરેક તબક્કે દરેક શાખાના આવાં બધાં
સૂત્રો પર આ નિયમો લાગુ કરી શકાય છે
(આ સ્વરૂપનું ન હોય એવા સૌથી ઉપરના ચિહ્નિત
!!{formula} પર કામ કર્યા પછી). વાપરવાનું પદ
$t$ એ શાખામાં પહેલેથી આવતાં !!{constant}s
(અથવા કોઈ !!{constant} ન આવતો હોય તો
$\Obj c_0$) અને શરૂઆતનાં !!{formula}sમાં
આવતાં !!{function}s વડે બની શકે એવું પ્રથમ
પદ~$t$ છે, જેના માટે સંબંધિત ચિહ્નિત
!!{formula} $\sFmla{\True}{!B(t)}$ અથવા
$\sFmla{\False}{!B(t)}$ શાખામાં અગાઉ આવતું નથી.

આ સાબિતી-શોધથી બંધ કોષ્ટક-વૃક્ષ ન બને, તો
કાં તો તેમાં એવી સીમિત ખુલ્લી શાખા હોય છે
જેના બધા બિનઆણ્વિક !!{formula}s પર ખરાની
નિશાની છે, અથવા શોધ અનંત પણ સીમિત-શાખાવાળું
વૃક્ષ બનાવે છે, જેમાં અનંત શાખા હોવી જ પડે.
\olref[cpl]{sec}ની જેમ, આપણે એવી
!!a{structure}~$\Struct{M}$ વ્યાખ્યાયિત કરી
શકીએ કે શાખામાં $\sFmla{\True}{!A}$ આવે તો
$\Sat{M}{!A}$, અને $\sFmla{\False}{!A}$
આવે તો $\Sat/{M}{!A}$.
($\Theta =
\Setabs{!A}{\sFmla{\True}{!A} \text{ શાખા પર}}$
અને $\Xi =
\Setabs{!A}{\sFmla{\False}{!A} \text{ શાખા પર}}$ લો.)

આ રીતે શોધીને બનાવેલાં કોષ્ટક-વૃક્ષના દરેક
શિરોબિંદુને સિક્વન્ટ સોંપવાથી તેને સરળતાથી
\Log{G3c}ની સાબિતીમાં રૂપાંતરિત કરી શકાય છે.
શરૂઆતનાં સૂત્રો $\Gamma \Sequent \Delta$
સિક્વન્ટને અનુરૂપ હોય, તો દરેક શરૂઆતના
સૂત્રને $\Gamma \Sequent \Delta$ ચિહ્ન આપો.
$\sFmla{\True}{!A}$ ચિહ્નિત !!{formula}
પર શાખા-વિસ્તારનો નિયમ લાગુ કરવાથી શાખા
વધે, તો પાનને $!A, \Pi \Sequent \Lambda$
ચિહ્ન આપો; ચિહ્નિત !!{formula}
$\sFmla{\False}{!A}$ પર તે
લાગુ થાય, તો પાનને $\Pi \Sequent \Lambda, !A$
ચિહ્ન આપો. કોષ્ટક-વૃક્ષમાં
\TRule{\True}{\star} નિયમ લાગુ થાય, ત્યાં
સિક્વન્ટ પર \LeftR{\star} નિયમ લગાવો
($!A$ને મુખ્ય સૂત્ર રાખીને); અને
\TRule{\False}{\star} લાગુ થાય, ત્યાં
\RightR{\star} નિયમ લગાવો. નિયમના
આધારવિધાનો કોષ્ટક-વૃક્ષમાં ઉમેરાયેલા
સંબંધિત શિરોબિંદુઓને ચિહ્નિત કરતા સિક્વન્ટો
છે. $\sFmla{\True}{!A}$ અને
$\sFmla{\False}{!A}$ વડે કોષ્ટક-વૃક્ષ બંધ
થાય, ત્યારે અંતિમ પાનનું ચિહ્ન
$!A, \Pi \Sequent \Lambda, !A$ સ્વરૂપનો
સિક્વન્ટ હશે. કોષ્ટક-વૃક્ષમાં ક્રમશઃ આવતા
કેટલાંક શિરોબિંદુઓને એક જ સિક્વન્ટ સોંપાશે;
પુનરાવર્તિત નકલો કાઢી નાખો. દાખલા તરીકે,
ઉપરનું કોષ્ટક-વૃક્ષ આ !!{proof}માં રૂપાંતરિત થાય છે:
\[\small
\Axiom$!A \fCenter !A, !B$
\RightLabel{\RightR{\lor}}
\UnaryInf$!A \fCenter !A \lor !B$
\Axiom$!A \fCenter !A, !C$
\RightLabel{\RightR{\lor}}
\UnaryInf$!A \fCenter !A \lor !C$
\RightLabel{\RightR{\land}}
\BinaryInf$!A \fCenter (!A \lor !B) \land (!A \lor !C)$
\Axiom$!B, !C \fCenter !A, !B$
\RightLabel{\RightR{\lor}}
\UnaryInf$!B, !C \fCenter !A \lor !B$
\RightLabel{\LeftR{\land}}
\UnaryInf$!B \land !C \fCenter !A \lor !B$
\Axiom$!B, !C \fCenter !A, !C$
\RightLabel{\RightR{\lor}}
\UnaryInf$!B, !C \fCenter !A \lor !C$
\RightLabel{\LeftR{\land}}
\UnaryInf$!B \land !C \fCenter !A \lor !C$
\RightLabel{\RightR{\land}}
\BinaryInf$!B \land !C \fCenter (!A \lor !B) \land
(!A \lor !C)$
\RightLabel{\LeftR{\lor}}
\BinaryInf$!A \lor (!B \land !C) \fCenter (!A \lor !B) \land
(!A \lor !C)$
\DisplayProof
\]

\end{document}
